\documentclass[11pt]{article}

\usepackage[margin=1in]{geometry}
\usepackage{amsmath,amssymb,amsthm}
\usepackage{enumitem}
\usepackage{hyperref}
\usepackage{xcolor}
\usepackage{booktabs}

\hypersetup{colorlinks=true, linkcolor=blue, urlcolor=blue, citecolor=blue}

\title{TOAE209/TOAH209 -- Design and Analysis of Algorithms\\[4pt]
\large Week 15: Theoretical Questions}
\author{Foreign Trade University}
\date{}

\begin{document}
\maketitle

\section*{Instructions}

Part A contains 18 questions requiring written explanations or short proofs. Part B is a
10-question quiz (true/false and multiple choice). Attempt every question on your own
before consulting Part C (Solutions) at the end of this document.

\section*{Part A: Theory Questions}

\begin{enumerate}[label=\textbf{A\arabic*.}, leftmargin=2.2em]

\item Define a \emph{randomized algorithm}. Explain the difference between a
\textbf{Las Vegas} algorithm and a \textbf{Monte Carlo} algorithm in terms of what is
guaranteed (correctness) and what is random (running time). Give one example of each from
this week.

\item State \emph{linearity of expectation} precisely. Why is it so useful that it holds
\textbf{even when the random variables are not independent}? Illustrate with the phrase
``decompose into indicators, then sum.''

\item Define an \emph{indicator random variable} $\mathbf{1}_A$ for an event $A$. Prove
that $\mathbb{E}[\mathbf{1}_A] = \Pr[A]$, and explain why this identity is the bridge that
lets us compute probabilities by computing expectations (and vice versa).

\item State the \emph{union bound} and prove it for two events from the inclusion-exclusion
identity $\Pr[A\cup B] = \Pr[A] + \Pr[B] - \Pr[A\cap B]$. Why does the bound require no
independence assumption?

\item In the contention-resolution model, $n$ identical processes each attempt to access a
shared resource with probability $p$ per round; an attempt succeeds only if exactly one
process attempts that round. Derive the probability that a fixed process $P_i$ succeeds in
a given round, and show it is maximized at $p = 1/n$.

\item Continuing Q5 with $p = 1/n$: show that $P_i$ succeeds in a given round with
probability at least $\frac{1}{en}$, using $(1 - 1/n)^{n-1} \ge 1/e$. Then bound the
probability that $P_i$ has \emph{not} succeeded after $t$ rounds.

\item Using the bound from Q6 and the union bound, prove that after $t = \lceil 2en\ln
n\rceil$ rounds, all $n$ processes have succeeded with probability at least $1 - 1/n$.
Identify exactly where the union bound is applied.

\item State randomized QuickSort. Explain why the analysis counts \emph{comparisons}, and
state the key structural fact that any two elements are compared \emph{at most once} over
the whole execution. Why is that fact true?

\item Let $z_1 < z_2 < \cdots < z_n$ be the sorted order. Prove that in randomized
QuickSort, $\Pr[z_i \text{ and } z_j \text{ are compared}] = \frac{2}{j - i + 1}$ for
$i < j$. (Hint: consider the first element of $\{z_i, \ldots, z_j\}$ chosen as a pivot.)

\item Using Q9 and linearity of expectation, prove that the expected number of comparisons
in randomized QuickSort is at most $2n\ln n = O(n\log n)$. Where does the harmonic-series
bound enter?

\item State randomized Quickselect for finding the $k$-th smallest element. Sketch the
argument that its expected running time is $O(n)$ (define a ``good'' pivot, bound the
expected number of pivots per phase, and sum the geometric series of shrinking subarray
sizes).

\item Describe Karger's contraction algorithm for global min-cut: what a \emph{contraction}
does to an edge, why \emph{parallel edges are kept but self-loops deleted}, and when the
algorithm stops.

\item In Karger's algorithm, fix a particular min-cut $C^*$ of size $k$. Prove that every
vertex has degree $\ge k$, and conclude that at any stage with $m$ super-vertices there are
at least $mk/2$ edges. Why does this matter for bounding the chance of contracting a
$C^*$-edge?

\item Prove that a single run of Karger's algorithm returns a fixed min-cut $C^*$ with
probability at least $1/\binom{n}{2}$. Show the telescoping product explicitly.

\item Karger's single-run success probability is only $\approx 2/n^2$. Explain
\emph{amplification}: how many independent repetitions $T$ suffice to find the min-cut with
probability $\ge 1 - 1/n$, and prove the bound using $1 - x \le e^{-x}$.

\item State the reservoir-sampling algorithm for selecting $k$ items from a stream of
unknown length. Prove that every item ends up in the reservoir with probability exactly
$k/n$ after $n$ items are seen. Point out where the product telescopes.

\item Define a \emph{universal} family of hash functions. For the family $h_{a,b}(x) =
((ax+b)\bmod p)\bmod m$, use indicator variables and linearity of expectation to show that
the expected number of colliding pairs when hashing $n$ keys into $m$ buckets is at most
$\binom{n}{2}/m \approx n^2/(2m)$.

\item \emph{Discussion (open-ended).} The Miller--Rabin primality test is Monte Carlo: it
never errs on a prime, but a composite may ``fool'' a random witness with probability
$\le 1/4$. Explain how to (a) drive the error below any target $\delta$ by repetition, and
(b) why a Monte Carlo algorithm with a \emph{checkable} answer can be converted into a Las
Vegas algorithm. Could primality testing be made Las Vegas this way? Discuss.

\end{enumerate}

\newpage
\section*{Part B: Quiz}

\begin{enumerate}[label=\textbf{B\arabic*.}, leftmargin=2.2em]

\item \textbf{(True/False)} A Las Vegas algorithm may return an incorrect answer but always
runs within a fixed time bound.

\item \textbf{(Multiple choice)} Linearity of expectation, $\mathbb{E}[\sum_i X_i] =
\sum_i \mathbb{E}[X_i]$, holds:
\begin{enumerate}[label=(\alph*)]
    \item Only if the $X_i$ are independent
    \item Only if the $X_i$ are identically distributed
    \item For any random variables, independent or not
    \item Only for indicator random variables
\end{enumerate}

\item \textbf{(Short answer)} In randomized QuickSort on $n$ distinct elements, what is the
probability that the $i$-th and $j$-th smallest elements ($i<j$) are ever compared?

\item \textbf{(True/False)} The expected number of comparisons made by randomized QuickSort
is $O(n \log n)$ on \emph{every} input, with no worst-case bad inputs.

\item \textbf{(Multiple choice)} Karger's contraction algorithm, on a graph with $n$
vertices, returns a fixed global min-cut in a single run with probability at least:
\begin{enumerate}[label=(\alph*)]
    \item $1/n$
    \item $1/\binom{n}{2}$
    \item $1/2$
    \item $1/n!$
\end{enumerate}

\item \textbf{(True/False)} Running Karger's algorithm $\Theta(n^2 \log n)$ times and
keeping the best cut finds the global min-cut with high probability.

\item \textbf{(Short answer)} In reservoir sampling of $k$ items from a stream, with what
probability is the $t$-th arriving item ($t > k$) accepted into the reservoir at the moment
it arrives?

\item \textbf{(Multiple choice)} The union bound states that $\Pr[\bigcup_i A_i]$ is:
\begin{enumerate}[label=(\alph*)]
    \item Exactly $\sum_i \Pr[A_i]$
    \item At most $\sum_i \Pr[A_i]$
    \item At least $\sum_i \Pr[A_i]$
    \item Equal to $\prod_i \Pr[A_i]$
\end{enumerate}

\item \textbf{(Short answer)} For a universal hash family hashing $n$ keys into $m$
buckets, approximately how many colliding pairs are expected?

\item \textbf{(True/False)} The Miller--Rabin primality test may report a prime as
``composite.''

\end{enumerate}

\newpage
\section*{Part C: Solutions}

\subsection*{Part A}

\begin{enumerate}[label=\textbf{A\arabic*.}, leftmargin=2.2em]

\item A \textbf{randomized algorithm} may consult a source of random bits (a random number
generator) during execution, so its behaviour on a fixed input can vary from run to run. A
\textbf{Las Vegas} algorithm is \emph{always correct}; only its running time is random (we
analyze \emph{expected} time) -- e.g.\ randomized QuickSort and Quickselect. A
\textbf{Monte Carlo} algorithm runs in bounded time but may be \emph{wrong with bounded
probability} (we analyze success probability, then amplify) -- e.g.\ Karger's min-cut and
Miller--Rabin. Slogan: Las Vegas gambles with time, Monte Carlo gambles with correctness.

\item \textbf{Linearity of expectation:} for any random variables $X_1,\ldots,X_n$ and
constants $a_i$, $\mathbb{E}[\sum_i a_i X_i] = \sum_i a_i \mathbb{E}[X_i]$. It holds
\emph{without any independence assumption} because expectation is just a weighted sum over
outcomes, and summation can be reordered regardless of how the variables correlate. This is
exactly what makes it powerful: to compute the expectation of a complicated quantity, we
\emph{decompose it into a sum of indicator variables} (one per ``elementary event''),
compute each easy probability $\mathbb{E}[\mathbf{1}_A] = \Pr[A]$, and sum -- never needing
to understand how the events interact.

\item $\mathbf{1}_A$ equals $1$ when $A$ occurs and $0$ otherwise. Then $\mathbb{E}[
\mathbf{1}_A] = 1\cdot\Pr[A] + 0\cdot\Pr[\bar A] = \Pr[A]$. This identity is the bridge
between probability and expectation: it lets us compute a probability by setting up an
expectation, and (via linearity) compute the expectation of a sum of events as a sum of
their probabilities -- the engine behind the QuickSort, reservoir-sampling, and hashing
analyses.

\item \textbf{Union bound:} $\Pr[\bigcup_i A_i] \le \sum_i \Pr[A_i]$. For two events, from
inclusion-exclusion $\Pr[A\cup B] = \Pr[A] + \Pr[B] - \Pr[A\cap B] \le \Pr[A] + \Pr[B]$,
since $\Pr[A\cap B]\ge 0$. The general case follows by induction. No independence is needed
because we only ever \emph{drop} the (non-negative) intersection terms, which is valid for
any events.

\item Process $P_i$ succeeds in a round iff it attempts (probability $p$) and \emph{no
other} process attempts (each of the other $n-1$ stays silent, probability $(1-p)^{n-1}$),
so $\Pr[\text{success}] = p(1-p)^{n-1}$. Differentiating $f(p) = p(1-p)^{n-1}$ and setting
$f'(p) = (1-p)^{n-1} - p(n-1)(1-p)^{n-2} = 0$ gives $(1-p) = p(n-1)$, i.e.\ $1 = pn$, so
$p = 1/n$ maximizes the success probability.

\item With $p = 1/n$, $\Pr[\text{success in a round}] = \frac1n(1 - \frac1n)^{n-1} \ge
\frac1n\cdot\frac1e = \frac{1}{en}$, using $(1-1/n)^{n-1}\ge 1/e$. Since rounds are
independent, the probability $P_i$ fails every one of the first $t$ rounds is at most
$(1 - \frac{1}{en})^t \le e^{-t/(en)}$ (using $1 - x \le e^{-x}$).

\item Set $t = \lceil 2en\ln n\rceil$. Then $\Pr[P_i \text{ not yet succeeded}] \le
e^{-t/(en)} \le e^{-2\ln n} = 1/n^2$. The bad event ``some process has not succeeded'' is
the union of the $n$ events ``$P_i$ not yet succeeded.'' By the \textbf{union bound}, its
probability is at most $n\cdot 1/n^2 = 1/n$. Hence with probability $\ge 1 - 1/n$ all
processes have succeeded -- the union bound is applied exactly at the step combining the $n$
per-process failure probabilities.

\item \textbf{Randomized QuickSort:} pick a pivot uniformly at random, partition into
elements smaller and larger, recurse on each side, concatenate. We count \emph{comparisons}
because they dominate the running time (partitioning is comparison-driven). Key fact: two
elements are compared \emph{at most once}, because a comparison occurs only between a pivot
and a non-pivot, and once an element is chosen as a pivot it is removed from all recursive
subproblems, so it never gets compared to anyone again.

\item Consider $Z_{ij} = \{z_i, z_{i+1}, \ldots, z_j\}$, the $j-i+1$ elements between $z_i$
and $z_j$ inclusive. Track the \emph{first} element of $Z_{ij}$ chosen as a pivot. If it is
$z_i$ or $z_j$, then at that moment $z_i$ and $z_j$ are still in the same subarray, so the
pivot is compared to the other -- they \emph{are} compared. If it is any interior $z_k$
($i<k<j$), then $z_i$ and $z_j$ are split into different subarrays and never compared.
Because pivots are uniform, the first pivot from $Z_{ij}$ is equally likely to be any of
the $j-i+1$ elements; exactly $2$ of them ($z_i, z_j$) cause a comparison. Hence the
probability is $\frac{2}{j-i+1}$.

\item Let $X = \sum_{i<j} X_{ij}$ be the total number of comparisons, $X_{ij}$ the
indicator that $z_i, z_j$ are compared. By linearity, $\mathbb{E}[X] = \sum_{i<j}
\frac{2}{j-i+1}$. Substituting $k = j-i+1 \in \{2,\ldots,n\}$ and over-counting,
\[
\mathbb{E}[X] \le \sum_{i=1}^{n}\sum_{k=2}^{n}\frac{2}{k} = 2n\sum_{k=2}^{n}\frac1k \le 2n\ln n,
\]
where the final inequality is the \textbf{harmonic-series bound} $\sum_{k=1}^n 1/k \le
1 + \ln n$. So $\mathbb{E}[X] = O(n\log n)$.

\item \textbf{Quickselect:} pick a random pivot, partition, and recurse only into the
\emph{one} side containing the $k$-th element (or return the pivot if it is the $k$-th).
\emph{Expected $O(n)$:} call a pivot \emph{good} if it falls between the $25$th and $75$th
percentiles -- probability $\ge 1/2$. A good pivot shrinks the surviving subarray to
$\le 3/4$ of its size. Group steps into phases ending at each good pivot; the expected
number of pivots per phase is $\le 2$ (geometric, success prob.\ $\ge 1/2$), and sizes
shrink by a factor $\ge 3/4$ each phase. Total expected work $\le \sum_{j\ge 0} 2c
(3/4)^j n = 8cn = O(n)$, the geometric series converging.

\item A \textbf{contraction} of edge $(u,v)$ merges $u$ and $v$ into a single
super-vertex; all other edges incident to $u$ or $v$ now attach to the merged vertex.
\textbf{Parallel edges are kept} because the multiplicity of edges between super-vertices
is exactly what records how many original edges cross a potential cut -- discarding them
would corrupt the cut count. \textbf{Self-loops are deleted} because an edge whose both
endpoints have been merged into the same super-vertex can never cross any cut. The
algorithm \textbf{stops when only two super-vertices remain}; the edges between them form
the returned cut.

\item If some vertex $v$ had degree $< k$, then the cut separating $\{v\}$ from the rest
would have size $< k$, contradicting that the min-cut is $k$; so every vertex has degree
$\ge k$. At a stage with $m$ super-vertices, the sum of degrees is $\ge mk$, and since each
edge contributes to two degrees, the number of edges is $\ge mk/2$. This matters because it
\emph{lower-bounds the total number of edges}, which \emph{upper-bounds} the chance that a
uniformly random edge is one of the $k$ cut edges: $k / (mk/2) = 2/m$.

\item Fix a min-cut $C^*$ of size $k$. The algorithm returns $C^*$ iff it never contracts a
$C^*$-edge. At step $i$ there are $n-i+1$ super-vertices, hence $\ge (n-i+1)k/2$ edges, of
which $k$ are in $C^*$; so the chance of contracting a $C^*$-edge is $\le \frac{2}{n-i+1}$.
Thus
\[
\Pr[C^* \text{ survives}] \ge \prod_{i=1}^{n-2}\Bigl(1 - \frac{2}{n-i+1}\Bigr)
= \prod_{i=1}^{n-2}\frac{n-i-1}{n-i+1}
= \frac{n-2}{n}\cdot\frac{n-3}{n-1}\cdots\frac{2}{4}\cdot\frac{1}{3} = \frac{2}{n(n-1)} = \frac{1}{\binom{n}{2}},
\]
the product telescoping (each numerator cancels a denominator two steps later).

\item \textbf{Amplification:} run the algorithm $T$ times independently and keep the
smallest cut found. Each run misses $C^*$ with probability $\le 1 - 1/\binom{n}{2}$, so all
$T$ miss with probability $\le (1 - 1/\binom n2)^T \le e^{-T/\binom n2}$ (using $1-x \le
e^{-x}$). Choosing $T = \binom{n}{2}\ln n$ makes this $\le e^{-\ln n} = 1/n$. Hence with
probability $\ge 1 - 1/n$ some run finds $C^*$ and the best-of-$T$ answer is the min-cut.
This is $\Theta(n^2 \log n)$ runs.

\item \textbf{Reservoir sampling:} keep the first $k$ items; for each later item $x_t$
($t>k$), with probability $k/t$ replace a uniformly random reservoir element with $x_t$.
\emph{Claim: item $x_t$ ends in $R$ with probability $k/n$.} It is in $R$ right after step
$t$ with probability $k/t$ (for $t\le k$ this is $1$). For each later step $s>t$, $x_t$ is
evicted only if step $s$ replaces (prob.\ $k/s$) \emph{and} picks $x_t$'s slot (prob.\
$1/k$), i.e.\ with probability $1/s$; it survives with probability $(s-1)/s$. Therefore
\[
\Pr[x_t\in R] = \frac{k}{t}\prod_{s=t+1}^{n}\frac{s-1}{s} = \frac{k}{t}\cdot\frac{t}{n} = \frac{k}{n},
\]
the product telescoping. Since every item is equally likely ($k/n$), $R$ is a uniform
random $k$-subset.

\item A family $\mathcal H$ is \textbf{universal} if for all distinct $x \ne y$,
$\Pr_{h\in\mathcal H}[h(x)=h(y)] \le 1/m$. For the $n$ keys, let $X_{xy}$ indicate that the
pair $(x,y)$ collides. By linearity, the expected number of colliding pairs is
\[
\mathbb{E}\Bigl[\sum_{x<y} X_{xy}\Bigr] = \sum_{x<y}\Pr[h(x)=h(y)] \le \binom{n}{2}\cdot\frac1m = \frac{n(n-1)}{2m} \approx \frac{n^2}{2m}.
\]
The $(ax+b\bmod p)\bmod m$ family with random $a\in\{1,\ldots,p-1\}, b\in\{0,\ldots,p-1\}$
is universal, so this bound applies.

\item (a) Each independent Miller--Rabin witness exposes a composite with probability
$\ge 3/4$ (fails to expose with prob.\ $\le 1/4$); after $r$ independent witnesses the
probability a composite slips through is $\le (1/4)^r$, which is below any target $\delta$
for $r \ge \log_4(1/\delta)$. (b) A Monte Carlo algorithm whose answer can be
\emph{efficiently verified} can be made Las Vegas: repeat until the verifier accepts -- the
output is then always correct, and the (now random) running time is the expected number of
trials times the per-trial cost. For \emph{primality}, the situation is subtle: a single
``composite'' verdict is certain (a witness is a certificate of compositeness), but a
``prime'' verdict from Miller--Rabin is not directly checkable in the same cheap way.
Historically, deterministic polynomial-time primality (AKS, 2002) settled the decision
problem outright; randomized Las Vegas \emph{primality proving} (e.g.\ ECPP) also exists but
is more involved. So yes, primality can be decided/proved without one-sided error, just not
by trivially wrapping Miller--Rabin in a cheap verifier.

\end{enumerate}

\subsection*{Part B}

\begin{enumerate}[label=\textbf{B\arabic*.}, leftmargin=2.2em]

\item \textbf{False.} That describes a \emph{Monte Carlo} algorithm. A Las Vegas algorithm
is always correct; its \emph{running time} is the random quantity.

\item \textbf{(c) For any random variables, independent or not.} This is precisely what
makes linearity of expectation so useful.

\item $\dfrac{2}{j - i + 1}$ (probability that the first pivot drawn from $\{z_i,\ldots,
z_j\}$ is one of the two endpoints).

\item \textbf{True.} Random pivoting removes any worst-case input: the $O(n\log n)$ bound
is on the \emph{expected} comparisons for every fixed input.

\item \textbf{(b) $1/\binom{n}{2}$}, equivalently $\dfrac{2}{n(n-1)}$.

\item \textbf{True.} $T = \binom{n}{2}\ln n = \Theta(n^2\log n)$ repetitions drive the
failure probability to $\le 1/n$.

\item $k/t$ -- the $t$-th item ($t>k$) is accepted into the reservoir with probability
$k/t$ when it arrives.

\item \textbf{(b) At most $\sum_i \Pr[A_i]$.} The union bound is an upper bound; equality
holds only when the events are disjoint.

\item About $\dfrac{n^2}{2m}$ (precisely $\binom{n}{2}/m$) colliding pairs in expectation.

\item \textbf{False.} Miller--Rabin has \emph{one-sided} error: it never misclassifies a
prime (it always reports a prime as ``probably prime''); only a composite can be reported
as ``probably prime.''

\end{enumerate}

\end{document}
